% =====================================================================
% Resumen ejecutivo
% =====================================================================
\section*{Resumen ejecutivo}
\addcontentsline{toc}{section}{Resumen ejecutivo}

La Oficina de Investigación e Información de la UPME busca municipios PDET
candidatos para pilotos de energía solar. Este informe responde la pregunta
central del proyecto: \textbf{cuántas edificaciones hay y cuánta área de
techos tiene cada uno de los 170 municipios PDET}, según dos fuentes abiertas
de huellas de edificios, con una metodología reproducible construida sobre
MongoDB.

\textbf{Qué se hizo.} Se cargaron en MongoDB los límites oficiales de los 170
municipios (MGN 2025 del DANE) y las huellas de edificios de Google Open
Buildings v3 (2\,694\,786 edificios en PDET) y Microsoft Building Footprints
(1\,761\,048). Cada edificio se asignó a un único municipio por su centroide
con una consulta \campo{\$geoWithin} apoyada en un índice \campo{2dsphere}, y
las métricas por municipio se calcularon con un pipeline de agregación dentro
de la base. Todo el flujo, de la descarga a los mapas, corre con un solo
comando (\campo{make all}) y se reprodujo desde una base vacía.

\begin{table}[H]
  \centering\small
  \caption{Cifras principales (Google con confianza $\geq 0{,}65$).}
  \label{tab:resumen-ejecutivo}
  \begin{tabularx}{\textwidth}{@{}L r r@{}}
    \toprule
    & \textbf{Google} & \textbf{Microsoft} \\
    \midrule
    Municipios PDET (16 subregiones, 19 departamentos) & \multicolumn{2}{c}{170} \\
    Edificios asignados a un municipio & 2\,694\,785 & 1\,761\,048 \\
    Área total de techos & 222,2\,km\textsuperscript{2} & 221,1\,km\textsuperscript{2} \\
    Área media por edificio & 82\,m\textsuperscript{2} & 126\,m\textsuperscript{2} \\
    Municipios sin ningún edificio & 5 & 0 \\
    \midrule
    Área respaldada por las dos fuentes, $\sum \min(A_G, A_M)$ & \multicolumn{2}{c}{201,6\,km\textsuperscript{2} (20\,156\,ha)} \\
    Potencia instalable estimada (orden de magnitud) & \multicolumn{2}{c}{16\,GWp (rango 12 a 20\,GWp)} \\
    \bottomrule
  \end{tabularx}
\end{table}

\textbf{Hallazgos.}
\begin{itemize}
  \item Las dos fuentes coinciden en el total (222 y 221\,km\textsuperscript{2})
        y en el orden de los municipios (correlación de Spearman 0,94), aunque
        Google divide la misma superficie en 1,5 veces más huellas, más
        pequeñas. Por eso se compara y se ordena por \emph{área}, no por
        número de edificios.
  \item No coinciden en el Pacífico: en el Chocó y la costa nariñense Google
        detecta muy poco o nada (cinco municipios sin edificios). 18 de los
        23 municipios donde las fuentes están en fuerte desacuerdo son de esa
        región.
  \item El área de techos se concentra: los 15 primeros municipios suman el
        39\,\% del área conservadora de los 170.
  \item El resultado es estable: subir el umbral de confianza de Google a
        0,70 o 0,75 reduce el total, pero no cambia cuáles municipios
        encabezan la lista.
\end{itemize}

\textbf{Recomendaciones para la UPME.}
\begin{enumerate}
  \item Priorizar los pilotos en los municipios del top 15 por área
        conservadora (tabla~\ref{tab:ranking}), empezando por los que además
        tienen el área concentrada en un casco urbano: \textbf{Santa Marta},
        \textbf{Santander de Quilichao} y \textbf{Apartadó}.
  \item No descartar los municipios del Chocó y del Pacífico por su área
        baja: ahí la cifra depende de la fuente, y deben verificarse con
        imágenes recientes o en campo antes de decidir.
  \item Antes de dimensionar un piloto, refinar el área útil y cruzarla con
        la irradiación y la demanda local (sección~\ref{sec:recomendaciones}).
\end{enumerate}

\newpage
